% Peta konsep materi pecahan - pakgun.gezytech.web.id
% Kompilasi: pdflatex tikz-peta-konsep.tex
\documentclass[tikz,border=10pt]{standalone}
\usepackage{amsmath}
\usetikzlibrary{arrows.meta,positioning}

\begin{document}
\pagecolor{#171D2B}
\definecolor{aksen}{HTML}{F7B538}
\color{white}

\begin{tikzpicture}[
  konsep/.style={draw=aksen, thick, rounded corners=4pt, fill=#1!25,
                 text=white, minimum width=3.2cm, minimum height=1.0cm,
                 font=\small, align=center},
  hubungan/.style={-{Stealth[length=3mm]}, draw=aksen, thick},
  label/.style={midway, above, font=\footnotesize, text=aksen},
]

  \node[konsep=orange] (PECAHAN) at (0,0) {Pecahan};

  \node[konsep=blue] (SENILAI)   at (0,-2.5) {Pecahan\\Senilai};
  \node[konsep=teal] (DESIMAL)   at (-4.5,-2.5) {Pecahan\\Desimal};
  \node[konsep=teal] (CAMPURAN)  at (5.2,-2.5) {Pecahan\\Campuran};

  \node[konsep=purple] (OPERASI) at (0,-5.0) {Operasi\\Pecahan};
  \node[konsep=purple] (SAMADENGAN) at (-4.5,-5.0) {Penjumlahan\\dengan Penyebut Sama};
  \node[konsep=purple] (BEDAPENYEBUT) at (5.2,-5.0) {Penjumlahan\\dengan Penyebut Berbeda};

  \draw[hubungan] (PECAHAN) -- (SENILAI) node[label] {beragam bentuk};
  \draw[hubungan] (PECAHAN) -- (DESIMAL) node[label, sloped] {ditulis dengan koma};
  \draw[hubungan] (PECAHAN) -- (CAMPURAN) node[label, pos=0.4, sloped, below] {bentuk campuran};
  \draw[hubungan] (SENILAI) -- (OPERASI) node[label, right=2pt] {dasar operasi};
  \draw[hubungan] (DESIMAL) -- (SAMADENGAN) node[label] {dicoba di penjumlahan};
  \draw[hubungan] (CAMPURAN) -- (BEDAPENYEBUT) node[label] {dicoba di penjumlahan};

\end{tikzpicture}
\end{document}
